\section{What Judges Use to Attribute Text}
\label{sec:mechanism}
\label{sec:watermark}

Most cases do not pass both controls, and those that pass are all on the
chat-app essays. We consider five possible explanations for how the judges
attribute texts:
\begin{itemize}
\item the texts contain too little information about their authors;
\item judges claim texts similar to what they would have written;
\item judges claim texts that are likely under their own model, the familiarity
that \citet{wataoka2024perplexity} relate to self-preference;
\item judges claim the texts they consider best
\citep{davidson2024anonymized, stamand2026sgtr};
\item judges detect a watermark added by their developer.
\end{itemize}

\paragraph{Information in the texts}
A random forest over 18 style features attributes \EngineAAcc{} of the chat-app
essays (95\% CI \EngineACI{}) and \OWClfPooled{} of the open-weight texts, more
than any judge. To test whether the judges can attribute an essay from its
vocabulary alone, we repeated the lineup with the words of each essay
shuffled. Claude's provider refused all 38 of Claude's shuffled lineups, and
Gemini returned no usable answer for \FSGeminiMissing{} of its 38
(Appendix~\ref{app:data}), so this condition covers four judges and omits the
judge with the largest self-advantage on the intact essays. On the \FSNLineups{} answered
lineups the four judges attribute \FSAcc{} of the essays correctly, compared
with \CXIntactAccFour{} for the same judges on the intact essays and 20\% by
chance (Figure~\ref{fig:shuffle} and Table~\ref{tab:fshuffle} in
Appendix~\ref{app:more}), and no self-advantage or discrimination is
significant. A bag-of-words classifier trained on labelled essays attributes
\BowShufAcc{} of them, consistent with \citet{sun2025stylometric}. Its input is
the same before and after shuffling, so this accuracy holds for both versions
by construction; it shows that vocabulary alone carries authorship information,
which the judges, who see no labelled examples, do not use. The condition
cannot show which cues the judges use on intact essays. Every reasoning trace
returned by Gemini and Grok describes the essays as scrambled
(Appendix~\ref{app:more}), so chance-level accuracy may reflect judges that
stop attributing when the text is visibly corrupted, and we do not interpret
it as evidence about word order.

\paragraph{Similarity to the judge's own answer}
For each frontier judge, we use its own answer to the same prompt as an
estimate of what it would have written, and test whether the judge names itself
more often on other models' texts that are similar to that answer. Across four
similarity measures (style features, word and character $n$-grams, and
function words), the pooled odds ratio per standard deviation of similarity is
between 0.89 and 1.04 in the lineup and between 0.91 and 1.09 one text at a
time, and every confidence interval includes 1. Of the 36 per-judge and pooled
tests, which are correlated, 1 reaches $p<0.05$. The per-judge tests use 13 to
79 self-attributions and cannot exclude odds ratios near 2. The same similarity
measures identify the true author in 74--79\% of cases, and the judges' answers
follow them only weakly. When judges are wrong in the lineup, the model they
name is the one whose answer is most similar to the text in 31.9\% of 630
errors, compared with 30.1\% under random reassignment, and one text at a time
the excess is similar (Appendix~\ref{app:more}). These excesses of 1 to 2
points are significant in a permutation test ($p\le 0.002$) that does not
account for the same essays recurring across judges.

\paragraph{Judge likelihoods}
For each open-weight judge and text, we compute the mean token log-probability
of the text under the judge. Conditioned on the original prompt, the judge's
own text is the most likely of the four candidates on
\OWLikCondRawMin{}--\OWLikCondRawMax{} of prompts (chance 25\%). This is
expected, since each text was sampled from its author. The judges did not see
the prompt in the lineup, one-text and yes/no formats; scored without it, the
judge's own text is the most likely on \OWLikUncRawMin{}--\OWLikUncRawMax{} of
prompts, with Qwen and Phi lowest. Subtracting the mean log-probability under
the other three judges, which no judge can compute, raises these shares to
\OWLikCondAccMin{} with the prompt and \OWLikUncAccMin{}--\OWLikUncAccMax{}
without it (Table~\ref{tab:owlik}). A judge's own probabilities therefore carry
authorship information, and this information is weaker when the prompt is
withheld. Among texts a judge did not write, we find no evidence that the texts
it claims have a higher relative likelihood than the texts it does not claim.
None of the \OWFamTests{} tests is significant after Holm correction, and the
only one with uncorrected $p<0.05$ (\OWFamMinWho{}; $p=$~\OWFamMinP{},
corrected \OWFamMinHolm{}) is in the opposite direction
(Appendix~\ref{app:more}). Five of the six tests in the one-text and yes/no
formats use 34 or fewer prompts and cannot exclude a moderate effect.

\paragraph{Text quality}
The four open-weight models rated every chat-app essay and every open-weight
text on a 1--9 scale. Frontier judges do not preferentially claim higher-rated
essays, and the lineup self-advantages of Claude, GPT and Grok remain
significant when quality is included as a covariate. Almost all frontier essays
receive ratings near the top of the scale, so this control is coarse. We
detect no self-preference in the open-weight judges. For each judge we take
the amount by which it rates its own texts above its peers' ratings of the
same texts, and subtract the same difference for texts by other authors. The result
is at most \OWQualSelfPrefMax{} points (Mistral, 95\% interval
\OWMistralQualSelfPrefCI{}) and at most \OWQwenQualSelfPref{} for the other
three judges. These ratings are compressed: \CXQualModalTop{} of the judges'
most likely ratings are 7, 8 or 9
(Appendix~\ref{app:openweight}), so a modest preference
could go undetected. No frontier judge rated any text, so we have no measure of
self-preference for the judges that pass both controls.

\paragraph{Watermarks}
The watermarking schemes we consider
\citep{kirchenbauer2023watermark, dathathri2024synthid} are detected with a
key or seeding function that the judges were not given. The chat-app essays
were written between 15 June and 31 July 2026, before Anthropic announced
watermarking of Claude text in August
\citep{anthropic2026watermark, anthropic2026marking}, so Claude's
self-advantage on these essays predates any announced Claude watermark. The
API texts were collected on 1 October 2026, after the announcement, so this
timing argument does not apply to them. On those texts no judge exceeds its
peer baseline in the lineup, and the other judges name Claude on its texts
almost as often as Claude does (\FXClaudeLineupPeer{} against
\FXClaudeLineupSelf{}), which a mark readable only by the author's model would
not produce. Text from the Gemini app is watermarked
\citep{deepmind2026synthid}, and Gemini does not recognise its own essays. The open-weight texts have no watermark, and the classifier attributes
\OWClfPooled{} of them. The shuffled lineup would also remove a watermark, but
because the judges are at chance in that condition, it does not distinguish
the two explanations.

\paragraph{Stated reasons}
Each frontier judgment includes a one-sentence reason, and we coded all 1,900
reasons (Appendix~\ref{app:rationales}). No reason refers to the judge itself
(for example, ``this is how I write''). Judges refer to every model, including
themselves, in the third person, and 95\% of one-text reasons cite a model's
typical style. The cited features are the same for correct and incorrect
self-attributions (Table~\ref{tab:reasons}): when Claude names itself one text at
a time, it cites hedging in 31 of 33 correct cases and in 12 of 13 incorrect
ones. Grok and Gemini attribute their own essays to Claude and describe Claude's
style as ``polished'' and ``elegant'', similar to the preference for frontier
model families reported by \citet{bai2025selfrecognition}. The reasons often
plausibly describe surface properties of the texts, but each judge applies the
same description of a model to all texts, so the reasons do not separate its
own texts from others'.
